\documentclass[11pt,a4paper]{article}
\usepackage[utf8]{inputenc}
\usepackage[T1]{fontenc}
\usepackage{lmodern}
\usepackage[french]{babel}
\usepackage{amsmath,amssymb,mathtools}
\usepackage{icomma}
\usepackage{xcolor}
\usepackage{tikz}
\usepackage{pgfplots}
\usepackage[most]{tcolorbox}
\usepackage{enumitem}
\usepackage{fancyhdr}
\usepackage[hidelinks]{hyperref}
\usepackage{titlesec}
\usepackage[a4paper,margin=2.2cm,headheight=26pt,headsep=10pt]{geometry}
\usetikzlibrary{arrows.meta,positioning,calc,intersections,angles,quotes}
\usepgfplotslibrary{fillbetween}
\pgfplotsset{compat=1.18}
\definecolor{primary}{RGB}{30,75,145}
\definecolor{primarydark}{RGB}{20,50,100}
\definecolor{defcol}{RGB}{0,114,178}
\definecolor{thmcol}{RGB}{0,135,110}
\definecolor{excol}{RGB}{0,130,85}
\definecolor{remarkcol}{RGB}{120,70,180}
\definecolor{attncol}{RGB}{200,40,55}
\definecolor{methcol}{RGB}{85,85,100}
\definecolor{areacol}{RGB}{120,160,230}
\definecolor{areacol2}{RGB}{230,150,80}
\newcommand{\dd}{\mathrm{d}}
\newcommand{\R}{\mathbb{R}}
\newcommand{\ua}{\,\text{u.a.}}
\newcommand{\tg}{\tan}\newcommand{\cotg}{\cot}
\tcbset{boxbase/.style={enhanced,breakable,boxrule=0.9pt,arc=2.5pt,left=3mm,right=3mm,top=2.3mm,bottom=2.3mm,fonttitle=\bfseries,coltitle=white,toptitle=0.6mm,bottomtitle=0.6mm}}
\newtcolorbox{methode}[1][]{boxbase,colback=methcol!7,colframe=methcol,sharp corners,title={\textsc{Comment faire}\ifx\relax#1\relax\relax\else\textmd{ : \itshape #1}\fi}}
\newtcolorbox{exemple}[1][]{boxbase,colback=excol!5,colframe=excol,title={\textsc{Exemple}\ifx\relax#1\relax\relax\else\textmd{ : \itshape #1}\fi}}
\newtcolorbox{idee}[1][]{boxbase,colback=defcol!5,colframe=defcol,title={\textsc{Idée clé}\ifx\relax#1\relax\relax\else\textmd{ : \itshape #1}\fi}}
\newtcolorbox{remarque}[1][]{boxbase,colback=remarkcol!6,colframe=remarkcol,title={\textsc{Remarque}\ifx\relax#1\relax\relax\else\textmd{ : \itshape #1}\fi}}
\newtcolorbox{attention}[1][]{boxbase,colback=attncol!6,colframe=attncol,title={\textsc{Attention}\ifx\relax#1\relax\relax\else\textmd{ : \itshape #1}\fi}}
\newtcolorbox{exo}[1][]{boxbase,colback={primary!4},colframe=primary,title={\textsc{Exercice #1}}}
\newtcolorbox{corr}[1][]{boxbase,colback={thmcol!5},colframe=thmcol,title={\textsc{Corrigé #1}}}
\tikzset{axe/.style={->,>=Stealth,black!65,line width=0.7pt},courbe/.style={primary,line width=1.3pt},courbe2/.style={areacol2!90!black,line width=1.3pt},dvert/.style={gray,dashed,line width=0.7pt}}
\pagestyle{fancy}\fancyhf{}
\fancyhead[L]{\small\textcolor{primarydark}{\textbf{Fiche 7} — Intégrale d'une fonction continue}}
\fancyhead[R]{\small\textcolor{primarydark}{\textsc{Thème 1 — Primitives \& règle du dU}}}
\fancyfoot[C]{\small\thepage}
\renewcommand{\headrulewidth}{0.8pt}
\titleformat{\section}{\large\bfseries\color{primarydark}}{\thesection}{0.6em}{}

\begin{document}

\section*{Niveau Facile — Corrigé détaillé}

\begin{corr}[1 — polynôme par linéarité]
Par linéarité, on intègre chaque terme séparément (forme 1 avec $U=x$) :
\[
I_1=3\int x^2\,\dd x+2\int x\,\dd x-5\int 1\,\dd x
=3\cdot\frac{x^3}{3}+2\cdot\frac{x^2}{2}-5x+C
=\boxed{\,x^3+x^2-5x+C.\,}
\]
\textbf{Vérification :} $(x^3+x^2-5x+C)'=3x^2+2x-5$. \checkmark
\end{corr}

\begin{corr}[2 — puissance de $x$]
Forme 1 avec $U=x$, $n=5$ ($n\neq -1$) :
\[
I_2=\int x^5\,\dd x=\frac{x^{5+1}}{5+1}+C=\boxed{\,\frac{x^6}{6}+C.\,}
\]
\textbf{Vérification :} $\Bigl[\tfrac{x^6}{6}\Bigr]'=\tfrac{6x^5}{6}=x^5$. \checkmark
\end{corr}

\begin{corr}[3 — la forme $\dd U/U$]
C'est le cas exclu de la forme puissance ($n=-1$). On utilise la forme 2 avec $U=x$ :
\[
I_3=\int \frac{\dd x}{x}=\boxed{\,\ln|x|+C.\,}
\]
\begin{attention}[valeur absolue obligatoire]
Écrire $\ln x+C$ serait faux sur $]-\infty,0[$. La forme correcte, valable sur tout intervalle
ne contenant pas $0$, est $\ln|x|+C$.
\end{attention}
\textbf{Vérification :} pour $x>0$, $(\ln x)'=\tfrac1x$ ; pour $x<0$, $(\ln(-x))'=\tfrac{-1}{-x}=\tfrac1x$. \checkmark
\end{corr}

\begin{corr}[4 — sinus et cosinus]
Forme 3, avec $U=x$ :
\[
\text{(a) } \int \cos x\,\dd x=\boxed{\,\sin x+C\,}\ ;\qquad
\text{(b) } \int \sin x\,\dd x=\boxed{\,-\cos x+C.\,}
\]
Le signe «\,$-$\,» de (b) vient de $(\cos x)'=-\sin x$ : c'est l'erreur classique à éviter.\\
\textbf{Vérification :} $(\sin x)'=\cos x$ \checkmark\quad et\quad $(-\cos x)'=\sin x$. \checkmark
\end{corr}

\begin{corr}[5 — exponentielle]
Forme 6 dans le cas $a=e$, avec $U=x$ (et $\ln e=1$) :
\[
I_5=\int e^{x}\,\dd x=\boxed{\,e^{x}+C.\,}
\]
\textbf{Vérification :} $(e^x)'=e^x$. \checkmark
\end{corr}

\end{document}
